% Livro A — capítulo, seção, texto, exemplo e exercício (prompt 03, §52).
\documentclass[11pt,oneside]{book}
% Preâmbulo comum das fixtures do scan (D48). Fontes Latin Modern com T1, para que a camada de
% texto traga os acentos como caracteres, e babel em português, para que \chapter imprima
% "Capítulo" como um livro brasileiro imprime.
\usepackage[T1]{fontenc}
\usepackage[utf8]{inputenc}
\usepackage[brazil]{babel}
\usepackage{lmodern}
\usepackage{amsmath,amssymb}
\usepackage{tikz}
\usepackage[a4paper,margin=2.5cm]{geometry}
\usepackage{enumitem}
\setlist[enumerate,2]{label=\alph*)}
\makeatletter\@ifundefined{c@chapter}{\newcounter{exemplo}}{\newcounter{exemplo}[chapter]}\makeatother
\newenvironment{exemplo}{\par\medskip\refstepcounter{exemplo}\noindent\textbf{Exemplo \theexemplo.}\ }{\par\medskip}
\newcommand{\blocoexercicios}[1][Exercícios]{\section*{#1}}

\begin{document}
\chapter{Funções}
\section{Funções lineares}
Uma função $f\colon \mathbb{R}\to\mathbb{R}$ é linear quando existe $a$ real tal que
$f(x) = ax$ para todo $x$. O gráfico de uma função linear é uma reta que passa pela origem.
Esta seção apresenta as propriedades básicas dessas funções e alguns exemplos.

\begin{exemplo}
A função $f(x) = 3x$ é linear, e $f(2) = 6$.

A verificação é direta: $f(x+y) = 3(x+y) = 3x + 3y = f(x) + f(y)$, e o mesmo vale
para o produto por escalar. Este segundo parágrafo é a resolução do exemplo, e
pertence a ele.
\end{exemplo}

\begin{center}
% TikZ e não `picture`: o `\line` do LaTeX desenha com glifos de uma fonte, e sairia do PDF como
% texto. O livro de verdade traz traço vetorial, e é isso que a fixture precisa ter.
\begin{tikzpicture}[x=1pt,y=1pt]
  \draw (10,10) -- (130,10);
  \draw (10,10) -- (10,80);
  \draw (10,80) -- (115,10);
  \node at (70,45) {$P$};
\end{tikzpicture}

Figura 1.1 --- O gráfico de uma função linear.
\end{center}

Toda função linear satisfaz $f(x+y) = f(x) + f(y)$, o que se verifica diretamente da definição.

\subsection{Domínio}
O domínio de uma função linear é o conjunto dos números reais.

\begin{exemplo}
Para $g(x) = -x$, temos $g(-4) = 4$.
\end{exemplo}

\blocoexercicios
\begin{enumerate}[label=\textbf{\arabic*.}]
\item Mostre que $f(x) = 5x$ é linear.
\item Calcule:
  \begin{enumerate}
  \item $f(1)$ para $f(x) = 2x$;
  \item $f(-3)$ para $f(x) = 7x$.
  \end{enumerate}
\item Explique por que $h(x) = x + 1$ não é linear.
\end{enumerate}

\section{Funções quadráticas}
Uma função quadrática tem a forma $f(x) = ax^2 + bx + c$, com $a \neq 0$.

\blocoexercicios
\begin{enumerate}[label=\textbf{\arabic*.}]
\item Encontre as raízes de $x^2 - 5x + 6$.
\item Determine o vértice da parábola $y = x^2 - 4x$.
\end{enumerate}
\end{document}
